\section{Computed two-leg values, solutions and relations}
\label{sec:at-computed}

Values of the two engine routes of Section \ref{sec:at-engines} (step2/strict/action/engine\_a, engine\_b, complete), equal entry by entry in $\mathbb Q(i)+\mathbb Q(i)\log2+\mathbb Q(i)\log3$, with the signs of Lemma \ref{lem:at-sformulas}. Units $u$; $L:=\log\tfrac32$; $T_0,\dots,T_4$, $g_1,\dots,g_4$, $s_0$, $S(\Delta)_d$, $S(P_b)_d$, $N$ as in Proposition \ref{prop:at-xi2} and Remark \ref{rem:at-twocomp}; $u_0:=\int(\rho(\AV)\eta,\lambda^*)$; $S(M_\omega)_{\rm spin}$ as in Proposition \ref{prop:at-two} (iii); $F^{\rm part}(\xi'):=u(L-\tfrac{79}{12})(T_1+T_2)(\xi')$ for $\xi'\in\Sigma$.

\begin{proposition}[Arity-one obstructions on two legs]
\label{prop:at-c-ob1}
\[
\begin{aligned}
&\hbar\Ob^{(0)}|_2=-\tfrac{27}8u\int(d^*dc,d^*\AV),\qquad \hbar\Ob^{(1)}(\Delta)|_2=\eqref{eq:at-reg2},\qquad \hbar\Ob^{(1)}(P_b)|_2=0,\qquad \hbar\Ob^{(1)}(R_r)|_2=0,\\
&\hbar\Ob^{(1)}(M_\omega)|_2=\tfrac13u\,Q\,S(M_\omega)_{\rm spin},\qquad \hbar\Ob^{(1)}(\xi)|_2=u\Big[\big(\tfrac{79}{24}-\tfrac12L\big)(g_1+g_3)+\tfrac{27}8g_2+0\cdot g_4\Big],\\
&\hbar\Ob^{(1)}(\eta)|_2=\sum_\mu x^\mu\,\hbar\Ob^{(1)}(\rho(e_\mu)\eta)|_2+u\Big[-\tfrac52\int(\eta,\lambda)\,d^*\AV-\tfrac52\int(\eta,\rho(d\AV)\lambda)+\big(-\tfrac32+3L\big)\int(\eta,\rho(d\Ap)\lambda)+\tfrac52\int(\rho(dc)\eta,\lambda^*)\Big]+Q(\cdot),
\end{aligned}
\]
the intrinsic part of the last line in the canonical density (Definition c11-canonical); along $K_b$ the obstruction is in the span of the canonical data with position degree $\le2$, with pole sum $\tfrac94\{\int(A,A),S(K_b)\}$. $\lc^{(1)}(\xi)|_2=\{uN,S_2(\xi)\}=2u(T_1+T_2-T_3-T_4)$, so $l_\xi=0$ in Remark \ref{rem:at-logexact} on two legs. The Wess--Zumino identities of Lemma \ref{lem:at-sformulas} (ii) hold on every entry; the printed form before the revision (opposite sign) fails on every entry where it was tested.
\end{proposition}

\begin{proposition}[Two-leg solutions with all free directions]
\label{prop:at-c-sol}
$F_1^{(0)}|_2=\tfrac{27}{16}u\int(d^*\AV,d^*\AV)$ plus the $s$-closed arity-zero directions, and
\[
\begin{aligned}
\hbar F_1^{(1)}(\Delta)|_2&=u\Big[\tfrac{22}3\int(A,A^*)-\tfrac{47}3\int(\lambda^*,\lambda)-\tfrac83\int(c^*,c)\Big]+u\sum_d\theta^\Delta_dS(\Delta)_d+\kappa_V\,uQ\!\int(\iota_x\AV,c^*)+\kappa_\lambda\,uQ\!\int(\rho(x)\lambda^*,\lambda^*),\\
\hbar F_1^{(1)}(P_b)|_2&=u\sum_d\theta^P_dS(P_b)_d+(\text{the }b\text{-parts of the }\kappa\text{ directions}),\\
\hbar F_1^{(1)}(\xi)|_2&=F^{\rm part}(\xi)+u\,\theta^\xi_VS^V_2(\xi)+u\,\theta^\xi_\perp S^\perp_2(\xi)+\{N_f,S_2(\xi)\}\text{-shifts},\\
\hbar F_1^{(1)}(\eta)|_2&=\sum_\mu x^\mu F^{\rm part}(\rho(e_\mu)\eta)-\big(3L-\tfrac32\big)u\,s_0-\tfrac52u\,u_0+u\,\theta^{\eta}_VS^V(\eta)+u\,\theta^{\eta}_\perp S^\perp(\eta),\\
\hbar F_1^{(1)}(M_\omega)|_2&=\tfrac13u\,S(M_\omega)_{\rm spin}+u\sum_d\theta^M_dS(M_\omega)_d,\qquad \hbar F_1^{(1)}(R_r)|_2=u\sum_d\theta^R_dS(R_r)_d,\\
\hbar F_1^{(1)}(K_b)|_2&=u\int(b,x)\Big\{-\tfrac{22}3(A,A^*)+16(\lambda,\lambda^*)+\tfrac83(c,c^*)\Big\}+\tfrac23u\int\big\{(\iota_bA)(\iota_xA^*)-(\iota_xA)(\iota_bA^*)\big\}\\
&\quad-\tfrac13u\int b_\mu x_\nu(\lambda,\rho_\mu\rho_\nu\lambda^*)+u\sum_d\theta^K_dS(K_b)_d+\sum_{j=1}^4\kappa^K_jG^K_j,
\end{aligned}
\]
with $Q$ of the first bracket equal to $\hbar\Ob^{(1)}(\Delta)|_2$ (gate G-F1Delta), $G^K_1,\dots,G^K_4$ the four second-order free-field symmetries of weight $3$ (closed, not exact, not of the form $S(z)$ on two legs; listed in action/complete/results.json, F1(K\_b)|2, free\_B), and the $x$-linear part of $\hbar F_1^{(1)}(K_b)|_2$ equal to $-(b,x)$ times the particular part along $\Delta$ plus the spin term of the rotation part of $d_{b'}v$ ($16-\tfrac13=\tfrac{47}3$). On two legs $H^0(Q)$ vanishes along $\Delta$, $P_b$, $M_\omega$, $R_r$, $\Sigma$, and is $1$-, $4$-dimensional along $\Sigma^*$ (position degree $\le1$), $K_b$ (position degree $\le2$); in both cases the residual has no component along it.
\end{proposition}

\begin{proposition}[Two-leg arity-two equations]
\label{prop:at-c-ar2}
On two legs, with the solutions of Proposition \ref{prop:at-c-sol}, the classes of the arity-two residuals vanish exactly under the following conditions (both routes):
\[
\begin{array}{l|l}
\text{pair} & \text{condition}\\
\hline
(\Delta,P_b) & \theta^\Delta_\perp=\theta^\Delta_\lambda=\theta^\Delta_V=-\tfrac13\quad(o_d=-\tfrac23u)\\
(\Delta,\xi_\pm) & \theta^\Delta_\perp=\theta^\Delta_\lambda=\theta^\Delta_V=-\tfrac13,\ \theta^\xi\ \text{free}\\
(\xi_+,\xi_-) & \theta^P_d=\theta^{\xi_+}_V+\theta^{\xi_-}_V-\tfrac{27}4-2L\ (\text{all }d),\quad \theta^{\xi_+}_V+\theta^{\xi_-}_V=\theta^{\xi_+}_\perp+\theta^{\xi_-}_\perp\\
(\xi_\pm,\xi'_\pm) & \theta^{\xi_\pm}_V=\theta^{\xi_\pm}_\perp\\
(\xi_\pm,P_b) & \theta^P_V=\theta^P_\perp=\theta^P_\lambda\\
(\xi_+,\eta_-) & \theta^\xi_V-\theta^\xi_\perp=\theta^\eta_V-\theta^\eta_\perp-3L+\tfrac32\\
(\xi_+,\eta_+) & \theta^R_\perp=\theta^R_\lambda=:\tau,\ \theta^\Delta_d=\theta^M_d=\tau-\tfrac13,\ \theta^\xi_V+\theta^\eta_V=\tau+\tfrac{77}{12}+2L,\ \theta^\xi_\perp+\theta^\eta_\perp=\tau+\tfrac{95}{12}-L
\end{array}
\]
The obstructions: $\hbar\Ob^{(2)}(\Delta,\xi)|_2=u\big[8(T_1+T_2)-\tfrac{26}3(T_3+T_4)+\log\tfrac43\,Q(\psi_1+\psi_2)\big]$, $\psi_1=\int x^\mu(\Gamma(\xi,\lambda)^\mu,c^*)$, $\psi_2=\int x^\mu(A^*_\mu,(\xi,\lambda^*))$, solved by $\hbar F_1^{(2)}(\Delta,\xi)|_2=\log\tfrac43\,u(\psi_1+\psi_2)$ plus closed directions; $\hbar\Ob^{(2)}(\xi,P_b)|_2$ is $\tfrac14\log\tfrac43$ times a $Q$-exact functional; $\hbar\Ob^{(2)}(\Delta,P_b)|_2$ has no position part and no pole; every two-leg arity-two pole sum vanishes. With $\theta^\xi_V=\theta^\xi_\perp=:t_\xi$, $\theta^\eta_V=:t_\eta$ (the position part of $F_1^{(1)}(\eta)$ is then $x^\mu F_1^{(1)}(\rho(e_\mu)\eta)$ with one uniform coefficient) and uniform $t_P$, $t_\Delta$, $t_M$, $t_R$:
\[
t_\Delta=-\tfrac13,\qquad t_M=t_\Delta,\qquad t_R=t_\Delta+\tfrac13,\qquad t_{\xi_+}+t_{\xi_-}-t_P=\tfrac{27}4+2L,\qquad t_{\xi_+}+t_{\eta_+}-t_\Delta=\tfrac{27}4+2L .
\]
\end{proposition}

\begin{proof}
Computed: the cycles of Theorem \ref{thm:at-classes} and Proposition \ref{prop:at-classes2} in the two routes, the cohomology on complete bases fixed by representation counts (for $(\xi,\eta)$ decided on the value vectors directly, the generated $(\lambda,\lambda^*)$ position basis lacking one structure per chirality), and the residual formulas of Definition \ref{def:at-ob}. The arity-two generating function for two odd generators is $I+cS(\xi)+c'S(\xi')-cc'S^{(2)}(\xi,\xi')$, the sign forced by the classical master equation (gate GA-S2/GB-S2: $QS^{(2)}_2+S_2(P_{\Gamma|_V})-\{S_2(\xi),S_2(\xi')\}=0$). The last line: insert $\theta^\Delta=-\tfrac13$ into the $(\xi_+,\eta_+)$ row, giving $\tau=0$.
\end{proof}

\begin{corollary}[Scheme-independent content on the action side]
\label{cor:at-c-inv}
The inner direction ($t_z\mapsto t_z+c\lambda_z$, $\lambda=-1,+1,-2,+2,0$ on $\Sigma,\Sigma^*,P,K$ and the even rotations and $\Delta$) and the outer direction ($t_{\xi_\pm}\mapsto t_{\xi_\pm}\pm a$, $t_{\eta_\pm}\mapsto t_{\eta_\pm}\mp a$) leave invariant exactly the combinations
\[
t_\Delta=-\tfrac13,\qquad t_M-t_\Delta=0,\qquad t_R-t_\Delta=\tfrac13,\qquad t_{\xi_+}+t_{\xi_-}-t_P=\tfrac{27}4+2L,\qquad t_{\xi_+}+t_{\eta_+}-t_\Delta=\tfrac{27}4+2L,
\]
and the sector uniformity of $t_P$, $t_\xi$, $t_\eta$; these are the one-loop parts of $[D,P]=-2P$, $[D,Q]=-Q$, $[Q,P]=0$, $[Q,Q]=P$ and $\{Q,S\}=D+M+R$ in the form of Theorem \ref{thm:at-algebra}. In a translation-covariant scheme $t_P=t_\Delta=-\tfrac13$, hence $t_{\xi_+}+t_{\xi_-}=\tfrac{77}{12}+2L$ and $t_{\eta_+}=t_{\xi_-}$. Scheme artefacts: the individual $t_{\xi_\pm}$, $t_{\eta_\pm}$, $t_P$, $t_K$, $\theta^K_d$, $\kappa^K_j$, $\kappa_V$, $\kappa_\lambda$, the particular parts with $\log\tfrac32$, and every representative above.
\end{corollary}

\begin{remark}[Two-component forms]
\label{rem:at-c-tc}
With Remark \ref{rem:at-twocomp}: $\int(A,A^*)=\int d^4x\{(A_m,A^*_m)+\tfrac14(A_{ij},A^{*ij})\}$, $\int(\lambda^*,\lambda)=\int(\lambda,\lambda^*)=\int d^4x\{(\lambda_\alpha{}^i,\lambda^{*\alpha}{}_i)-(\bar\lambda_{\dot\alpha i},\bar\lambda^{*\dot\alpha i})\}$, $\int(\iota_bA)(\iota_xA^*)=\int d^4x\,(b_mA_m,x_nA^*_n)$, and $S^{(2)}(\xi,\xi')=\int\{\tfrac14(\Gamma(\xi,\xi'),\Gamma(\lambda^*,\lambda^*))-\tfrac12(\xi,\lambda^*)(\xi',\lambda^*)-(\iota_{\Gamma(\xi,\xi')}A,c^*)\}$ with $\Gamma(\xi,\bar\xi')|_V^m=i\,\bar\xi'_{\dot\alpha i}\bar\sigma^{m\dot\alpha\alpha}\xi_\alpha{}^i$, $\Gamma|_\perp=0$ for $\xi\in S_+\otimes W$, $\bar\xi'\in S_-\otimes W^*$, and $\Gamma|_V=0$, $(\Gamma|_\perp)_{kl}=2\varepsilon_{jikl}\,\xi_\alpha{}^j\varepsilon^{\alpha\beta}\xi'_\beta{}^i$ for $\xi,\xi'\in S_+\otimes W$. On the Dirac column the $\rho_\mu\rho_\nu$ of the spin terms is $\gamma^\mu_W\gamma^\nu_W$ (Proposition c8-twocomp), on the undotted block $\sigma^\mu\bar\sigma^\nu$ and on the dotted block $\bar\sigma^\mu\sigma^\nu$.
\end{remark}

